\documentclass[11pt, letterpaper]{article}
\input{../../homework.tex}
\usepackage{titlepageBU}
\title{Problem Set 2}
\courseID{PY 451}
\courseSection{A1}
\begin{document}
\makeproblem
\section*{Problem 1}
	i feel bad for whoever's grading this
\begin{enumerate}
	\item 
		\[
			\bra{b_1} A \ket{a_1} =\begin{pmatrix}
			-1 & -i
			\end{pmatrix}^* \begin{pmatrix}
			1 & 2i \\
			-2 & i
			\end{pmatrix} \begin{pmatrix}
			3i \\
			2
			\end{pmatrix} = \begin{pmatrix}
			-1 & i
			\end{pmatrix} \begin{pmatrix}
			3i+4i \\
			-6i+2i
			\end{pmatrix} = -7i+4
		.\]
	\item 
		\[
			\bra{b_1} A^\dagger \ket{a_1} = \begin{pmatrix}
			-1 & -i
			\end{pmatrix}^* \begin{pmatrix}
			1 & -2 \\
			2i & i
			\end{pmatrix}^* \begin{pmatrix}
			3i \\
			2
			\end{pmatrix} = \begin{pmatrix}
			-1 & i
			\end{pmatrix} \begin{pmatrix}
			1 & -2 \\
			-2i & -i
			\end{pmatrix} \begin{pmatrix}
			3i \\
			2
			\end{pmatrix}=\begin{pmatrix}
			-1 & i
			\end{pmatrix} \begin{pmatrix}
			3i-4 \\
			-6-2i
			\end{pmatrix}
		\]
		\[
			 = (4-3i)+(2+6i)=6+3i
		.\]
	\item 
		\[
			\left( \bra{b_1} A \ket{a_1}  \right) ^{\dagger} =(-7i+4)^* = 4+7i
		.\]
	\item 
		\[
			A \ket{a_1} =\begin{pmatrix}
			1 & 2i \\
			-2 & i
			\end{pmatrix}\begin{pmatrix}
			3i \\
			2
			\end{pmatrix}=3i\begin{pmatrix}
			1 \\
			-2
			\end{pmatrix}+2 \begin{pmatrix}
			2i \\
			i
			\end{pmatrix}=\begin{pmatrix}
			7i \\
			-4i
			\end{pmatrix}
		.\]
	\item 
		\[
			\bra{b_1} A = \begin{pmatrix}
			-1 & -i
			\end{pmatrix}^* \begin{pmatrix}
			1 & 2i \\
			-2 & i
			\end{pmatrix} = \begin{pmatrix}
			-1 & i
			\end{pmatrix} \begin{pmatrix}
			1 & 2i \\
			-2 & i
			\end{pmatrix}=\begin{pmatrix}
			-1-2i & -2i-1
			\end{pmatrix}
		.\]
	\item 
		\[
			\braket{b_1}{a_1} =\begin{pmatrix}
			-1 & -i
			\end{pmatrix}^* \begin{pmatrix}
			3i \\
			2
			\end{pmatrix}=\begin{pmatrix}
			-1 & i
			\end{pmatrix}\begin{pmatrix}
			3i \\
			2
			\end{pmatrix}=-3i+2i=-i
		.\]
	\item
		\[
			\braket{a_1}{b_1} =\left( \braket{b_1}{a_1}  \right)^* = (-i)^*=i
		.\]
	\item 
		\[
			\ket{b_1} \bra{a_1} = \begin{pmatrix}
			-1 \\
			-i
			\end{pmatrix}\begin{pmatrix}
			3i & 2
			\end{pmatrix}^* = \begin{pmatrix}
			-1 \\
			-i
			\end{pmatrix}\begin{pmatrix}
			-3i & 2
			\end{pmatrix}=\begin{pmatrix}
			3i & -2 \\
			-3 & -2i
			\end{pmatrix}
		.\]
	\item
		\[
			A \ket{b_1} \bra{a_1} =\begin{pmatrix}
			1 & 2i \\
			-2 & i
			\end{pmatrix}\begin{pmatrix}
			3i & -2 \\
			-3 & -2i
			\end{pmatrix}=\begin{pmatrix}
			3i-6i & -2+4 \\
			-6i-3i & 4+2
			\end{pmatrix}=\begin{pmatrix}
			-3i & 2 \\
			-9i & 6
			\end{pmatrix}
		.\]
	\item
		\[
			A \ket{b_1} \bra{a_1} A= \begin{pmatrix}
			-3i & 2 \\
			-9i & 6
			\end{pmatrix}\begin{pmatrix}
			1 & 2i \\
			-2 & i
			\end{pmatrix}=\begin{pmatrix}
			-3i-4 & 6+2i \\
			-9i-12 & 18+6i
			\end{pmatrix}
		.\]
	\item why did we have to do this
		\[
			A^2 \ket{b_1} \bra{a_1} =\begin{pmatrix}
			1 & 2i \\
			-2 & i
			\end{pmatrix}\begin{pmatrix}
			-3i & 2 \\
			-9i & 6
			\end{pmatrix}=\begin{pmatrix}
			-3i-18 & 2+12i \\
			6i+9 & -4+6i
			\end{pmatrix}
		.\]
	\item
		\[
			\left( \ket{b_1} \bra{a_1}  \right) ^2 = \begin{pmatrix}
			3i & -2 \\
			-3 & -2i
			\end{pmatrix}\begin{pmatrix}
			3i & -2 \\
			-3 & -2i
			\end{pmatrix}=\begin{pmatrix}
			-9+6 & -6+4i \\
			-9i+6i & 6-4
			\end{pmatrix}=\begin{pmatrix}
			-3 & -2i \\
			-3i & 2
			\end{pmatrix}
		.\]
	\item
		\[
			A^{\dagger} \ket{a_1} =\begin{pmatrix}
			1 & -2 \\
			2i & i
		\end{pmatrix}^* \begin{pmatrix}
		3i \\
		2
		\end{pmatrix} = \begin{pmatrix}
		1 & -2 \\
		-2i & -i
	\end{pmatrix} \begin{pmatrix}
	3i \\
	2
	\end{pmatrix} = \begin{pmatrix}
	3i-4 \\
	6-2i
	\end{pmatrix}
		.\]
	\item why do $\ket{a_1} $ and $\ket{b_1} $ have subscripts if we're only using a single $a$ and $b$ like we could just do $\ket{a} $ and $\ket{b} $ and it would be the same
		\[
			\left( A \ket{b_1}  \right) ^{\dagger} =\left( \begin{pmatrix}
			1 & 2i \\
			-2 & i
			\end{pmatrix}\begin{pmatrix}
			-1 \\
			-i
			\end{pmatrix} \right) ^{\dagger} =\begin{pmatrix}
			-1+2 \\
			2+1
			\end{pmatrix}^{\dagger} =\begin{pmatrix}
			1 & 3
			\end{pmatrix}
		.\]
	\item
		\[
			A \ket{b_1} -\lambda \ket{b_1} =\begin{pmatrix}
			1 & 2i \\
			-2 & i
			\end{pmatrix}\begin{pmatrix}
			-1 \\
			-i
			\end{pmatrix}-(2+3i)\begin{pmatrix}
			-1 \\
			-i
			\end{pmatrix}=\begin{pmatrix}
			-1+2 \\
			2+1
			\end{pmatrix}-\begin{pmatrix}
			-2-3i \\
			-2i+3
			\end{pmatrix}
		\]
		\[
			=\begin{pmatrix}
			1 \\
			3
			\end{pmatrix}+\begin{pmatrix}
			2+3i \\
			-3+2i
			\end{pmatrix}=\begin{pmatrix}
			5+3i 
			2i
			\end{pmatrix}
		.\]
	\item
		\[
			\lambda A = (2+3i)\begin{pmatrix}
			1 & 2i \\
			-2 & i
			\end{pmatrix}=\begin{pmatrix}
			2+3i & 4i-6 \\
			-4-6i & 2i-3
			\end{pmatrix}
		.\]
	\item
		\[
			(\lambda -\braket{a_1}{b_1} )A=(2+3i-i)A=(2+2i)\begin{pmatrix}
			1 & 2i \\
			-2 & i
			\end{pmatrix}=4i-4\begin{pmatrix}
			2+2i &  \\
			-4-4i & 2i-2
			\end{pmatrix}
		.\]
	\item 
		\[
			\left( \braket{a_1}{a_1} +\braket{b_1}{b_1}  \right) ^{-1} =\left( \begin{pmatrix}
			3i & 2
			\end{pmatrix}^*\begin{pmatrix}
			3i \\
			2
			\end{pmatrix}+\begin{pmatrix}
			-1 & -i
			\end{pmatrix}^* \begin{pmatrix}
			-1 \\
			-i
			\end{pmatrix} \right)^{-1}
		\]
		\[
			 =\left( \begin{pmatrix}
			-3i & 2
			\end{pmatrix}\begin{pmatrix}
			3i \\
			2
			\end{pmatrix}+\begin{pmatrix}
			-1 & i
			\end{pmatrix}\begin{pmatrix}
			-1 \\
			-i
			\end{pmatrix} \right) ^{-1} =\left( 9+4+1+1 \right) ^{-1} =\frac{1}{15} 
		.\]
	\item 
		\[
			\left( \ket{a_1} \bra{a_1} +\ket{b_1} \bra{b_1}  \right) ^{-1} =\left( \begin{pmatrix}
			3i \\
			2
			\end{pmatrix}\begin{pmatrix}
			-3i & 2
			\end{pmatrix}+\begin{pmatrix}
			-1 \\
			-i
			\end{pmatrix}\begin{pmatrix}
			-1 & i
			\end{pmatrix} \right) ^{-1} =\left( \begin{pmatrix}
			9 & 6i \\
			-6i & 4
			\end{pmatrix}+\begin{pmatrix}
			1 & -i \\
			i & 1
			\end{pmatrix} \right) ^{-1} 
		\]
		\[
			=\begin{pmatrix}
				10 & 5i \\
			-5i & 5
			\end{pmatrix}^{-1} =\frac{1}{50+25} \begin{pmatrix}
			5 & -5i \\
			5i & 10
			\end{pmatrix}=\frac{1}{15} \begin{pmatrix}
			1 & -i \\
			i & 2
			\end{pmatrix}
		.\]
	\item 
		\[
			A A^{\dagger} =\begin{pmatrix}
			1 & 2i \\
			-2 & i
			\end{pmatrix}\begin{pmatrix}
			1 & -2 \\
			-2i & -i
			\end{pmatrix}=\begin{pmatrix}
			1+4 & -2+2 \\
			-2+2 & 4+1
			\end{pmatrix}=\begin{pmatrix}
			5 &  \\
			 & 5
			\end{pmatrix}
		.\]
	\item
		\[
			\left( A^2 \right) ^{\dagger} =\left( \begin{pmatrix}
			1 & 2i \\
			-2 & i
			\end{pmatrix}\begin{pmatrix}
			1 & 2i \\
			-2 & i
			\end{pmatrix} \right) ^{\dagger} =\begin{pmatrix}
			1-4i & 2i-2 \\
			-2-2i & -4i-1
			\end{pmatrix}^{\dagger} =\begin{pmatrix}
			1+4i & -2+2i \\
			-2-2i & -1+4i
			\end{pmatrix}
		.\]
	\item
		\[
			\left( A A^{\dagger}  \right) ^{\dagger} =\begin{pmatrix}
			5 &  \\
			 & 5
			\end{pmatrix}^{\dagger} =\begin{pmatrix}
			5 &  \\
			 & 5
			\end{pmatrix}
		.\]
\end{enumerate}
\section*{Problem 2}
\begin{enumerate}
	\item Since we know the state after the first measurement, we should know all subsequent probabilities. We have shown multiple times that measurements of a basis state in an orthogonal basis yield probabilities of $1/2$ for each state, so we have that every measurement that originated at $S_z \to \ket{+} $ should be the same. Using this fact, we find that the measurements for $S_z \to \ket{-} $ have 0.15 probability, to ensure the sum of the probabilities is 1.
		\[\begin{tikzcd}[row sep = 2.5em, column sep = 4em]
			&&&&0.1\\
			&&&&0.1\\
			&&&S_z\arrow[ru, "\ket{-}"']\arrow[ruu, "\ket{+}"]&0.1\\
			&&S_y\arrow[ru, "\ket{+}_y"]\arrow[r, "\ket{-}_y"']&S_x\arrow[ru, "\ket{+}_x"]\arrow[r, "\ket{-}_x"']&0.1\\
			\ket{\psi }\arrow[r]&S_z\arrow[ru, "\ket{+}"]\arrow[dr, "\ket{-} "'] \\
					    &&S_x\arrow[dr, "\ket{-}_x"']\arrow[r, "\ket{+}_x"]&S_y\arrow[dr, "\ket{-}_y"']\arrow[r, "\ket{+}_y"]&0.15\\
				   &&&S_z\arrow[ddr, "\ket{-}"']\arrow[dr, "\ket{+} "]&0.15\\
				   &&&&0.15\\
				   &&&&0.15
		\end{tikzcd}\]
		
	\item First we find the norm of the given vector.
		\[
			3\ket{+}_x+5i\ket{-}_x+7\ket{+}_y \cong \frac{3}{\sqrt{2} } \begin{pmatrix}
			1 \\
			1
			\end{pmatrix}+\frac{5i}{\sqrt{2} } \begin{pmatrix}
			1 \\
			-1
			\end{pmatrix}+\frac{7}{\sqrt{2} } \begin{pmatrix}
			1 \\
			i
			\end{pmatrix} = \frac{1}{\sqrt{2} } \begin{pmatrix}
			10+5i \\
			3+2i
			\end{pmatrix}
		.\]
		\[
			\left\| \ket{\psi }  \right\| = \frac{1}{\sqrt{2} } \sqrt{\left| 10+5i \right|^2+\left| 3+2i \right|^2} =\frac{1}{\sqrt{2} } \sqrt{125+13} =\frac{\sqrt{138} }{\sqrt{2} } =\sqrt{69} 
		.\]
	It follows that
	\[
		\ket{\psi}  = \frac{1}{\sqrt{69} } \left( 3\ket{+}_x+5i\ket{-}_x+7\ket{+}_y \right) 
	.\]
	\item First, we compute the probabilities after the Z Stern-Gerlach machine. Note that
		\[
			\ket{\psi }= \frac{10+5i}{\sqrt{138} } \ket{+}+\frac{3+2i}{\sqrt{138} } \ket{-}
		.\]
		We have
		\begin{align*}
			\mathcal{P}_+&= \left| \bra{+}P_+ \ket{\psi } \right| ^2\\
				     &=\left| \braket{+}{+}\braket{+}{\psi } \right| ^2\\
				     &=\left| \braket{+}{\psi } \right| ^2\\
				     &=\left| \frac{1}{\sqrt{138} } \left( (10+5i)\braket{+}{+}+(3+2i)\braket{+}{-} \right)  \right| ^2\\
				     &=\left| \frac{10+5i}{\sqrt{138} }  \right| ^2\\
				     &=\frac{125}{138} 
		.\end{align*}
		
		After the first measurement, the distribution (inside only the top half or bottom half) will be even, since all subsequent measurements have 50/50 probabilities that we have proven enough times in class that I don't feel I need to do it again here. This means to get each final probability, just multiply the denominator by 4 to obtain 125/552. The tree is on the next page because it was too big to fit on this one. Also we clearly have $\mathcal{P}_-=13/138$ to make the probabilities add to 1, so the other probabilities will be $13/552$.\newpage
		\[\begin{tikzcd}[row sep = 2.5em, column sep = 4em]
			&&&&125/552\\
			&&&&125/552\\
			&&&S_y\arrow[ru, "\ket{-}_y"']\arrow[ruu, "\ket{+}_y"]&125/552\\
			&&S_x\arrow[ru, "\ket{+}_x"]\arrow[r, "\ket{-}_x"']&S_y\arrow[ru, "\ket{+}_x"]\arrow[r, "\ket{-}_x"']&125/552\\
			\ket{\psi }\arrow[r]&S_z\arrow[ru, "\ket{+}"]\arrow[dr, "\ket{-} "'] \\
					    &&S_x\arrow[dr, "\ket{-}_x"']\arrow[r, "\ket{+}_x"]&S_y\arrow[dr, "\ket{-}_y"']\arrow[r, "\ket{+}_y"]&3/552\\
				   &&&S_y\arrow[ddr, "\ket{-}_y"']\arrow[dr, "\ket{+}_y"]&3/552\\
				   &&&&3/552\\
				   &&&&3/552
		\end{tikzcd}\]
\end{enumerate}
\section*{Problem 3}
Given a unit vector of the form $\hat{n} =(\sin \theta \cos \phi ,\sin \theta \sin \phi ,\cos \theta )$, the spin matrix $S_{\hat{n} } $ is given as
\[
	S_{\hat{n} } =\frac{\hbar}{2} \begin{pmatrix}
	\cos \theta  & \sin \theta e^{-i\phi }  \\
	\sin \theta e^{i\phi }  & -\cos \theta 
	\end{pmatrix}
.\]
In the problem, $\vec{n} =(\sin \theta ,0,\cos \theta )$ is of this form given that $\phi =0$. We thus have
\[
	S_{\vec{n} } =\frac{\hbar}{2}\begin{pmatrix}
	\cos \theta  & \sin \theta  \\
	\sin \theta  & -\cos \theta 
	\end{pmatrix}
.\]
We also have
\begin{align*}
	S_{\vec{n} } \ket{\psi_1} &\cong\frac{\hbar}{2} \begin{pmatrix}
	\cos \theta  & \sin \theta  \\
	\sin \theta  & -\cos \theta 
	\end{pmatrix}\begin{pmatrix}
	\cos (\theta /2) \\
	\sin (\theta /2)
	\end{pmatrix}\\
				  &=\frac{\hbar}{2} \begin{pmatrix}
				  \cos \theta \cos (\theta /2)+\sin \theta \sin (\theta /2) \\
				  \sin \theta \cos (\theta /2)-\cos \theta\sin (\theta /2)
				  \end{pmatrix}\\
				  &=\frac{\hbar}{2}\begin{pmatrix}
				  \cos \theta \cos (\theta /2)+2 \sin^2 (\theta /2)\cos (\theta /2) \\
				  2\sin( \theta/2) \cos^2 (\theta /2)-\cos\theta \sin (\theta /2)
				  \end{pmatrix}\\
				  &=\frac{\hbar}{2}\begin{pmatrix}
					  \left( 2 \cos ^2(\theta /2)-1 \right) \cos (\theta /2)+2 \sin ^2(\theta /2)\cos (\theta /2) \\
				  2 \sin (\theta /2)\cos ^2(\theta /2)-\left( 2 \cos ^2(\theta /2)-1 \right) \sin (\theta /2)
				  \end{pmatrix}\\
				  &=\frac{\hbar}{2}\begin{pmatrix}
				  2\left( \cos ^2(\theta /2) +\sin ^2(\theta /2)\right) \cos (\theta /2)-\cos (\theta /2) \\
				  2 \sin (\theta /2)\cos ^2(\theta /2)-\left(2 \cos ^2(\theta /2)\sin (\theta /2)-\sin (\theta /2)\right)
				  \end{pmatrix}\\
				  &=\frac{\hbar}{2}\begin{pmatrix}
				  \cos (\theta /2) \\
				  2\left( \cos ^2(\theta /2)-\cos ^2(\theta /2) \right) \sin (\theta /2)+\sin (\theta /2)
				  \end{pmatrix}\\
				  &=\frac{\hbar}{2} \begin{pmatrix}
				  \cos (\theta /2) \\
				  \sin (\theta /2)
				  \end{pmatrix}\\
				  &\cong \frac{\hbar}{2}\ket{\psi_1}
.\end{align*}
\begin{align*}
	S_{\vec{n} } \ket{\psi _2} &\cong \frac{\hbar}{2}\begin{pmatrix}
	\cos \theta  & \sin \theta  \\
	\sin \theta  & -\cos \theta 
	\end{pmatrix}\begin{pmatrix}
	\sin (\theta  /2) \\
	-\cos (\theta /2)
	\end{pmatrix}\\
				   &=\frac{\hbar}{2}\begin{pmatrix}
				   \cos \theta \sin (\theta /2)-\sin \theta \cos (\theta /2) \\
				   \sin \theta \sin (\theta /2)+\cos \theta \cos (\theta /2)
				   \end{pmatrix}\\
				   &=\frac{\hbar}{2}\begin{pmatrix}
				   -(\sin \theta \cos (\theta /2)-\cos \theta \sin (\theta /2) \\
				   \sin \theta \sin (\theta /2)+\cos \theta \cos (\theta /2)
				   \end{pmatrix}\\
				   &=\frac{\hbar}{2}\begin{pmatrix}
				   -\sin (\theta /2) \\
				   \cos (\theta /2)
				   \end{pmatrix}\\
				   &\cong -\frac{\hbar}{2}\ket{\psi _2} 
.\end{align*}
Here I've used the fact that the top entry of $S_{\vec{n} } \ket{\psi_2}  $ is the negative of the bottom entry of $S_{\vec{n} } \ket{\psi _1} $, so we can just plug in the answer from there, and similarly the bottom entry of $S_{\vec{n} } \ket{\psi _2} $ is the top entry of $S_{\vec{n} } \ket{\psi _1} $.

Consider $S_{\vec{n} }^2$.
\[
	\frac{\hbar}{2} \begin{pmatrix}
	\cos \theta  & \sin \theta  \\
	\sin \theta  & -\cos \theta 
	\end{pmatrix}\frac{\hbar}{2}\begin{pmatrix}
	\cos \theta  & \sin \theta  \\
	\sin \theta  & -\cos \theta 
	\end{pmatrix}=\frac{\hbar^2}{4}\begin{pmatrix}
	\cos ^2\theta +\sin ^2\theta  & \cos \theta \sin \theta -\sin \theta \cos \theta  \\
	\sin \theta \cos \theta -\cos \theta \sin \theta  & \sin ^2\theta +\cos ^2\theta 
	\end{pmatrix}= \frac{\hbar^2}{4} \mathbbm{1}
.\]
Therefore, $S_{\vec{n} }^2=\left(  \hbar^2/4 \right) \mathbbm{1}  $.
\end{document}
